% Παράρτημα B: Αναλυτικοί Πίνακες Αποτελεσμάτων

\chapter{Αναλυτικοί Πίνακες Αποτελεσμάτων}
\label{app:detailed-results}

\section{Επισκόπηση}

Το παρόν παράρτημα παρουσιάζει πλήρη αριθμητικά αποτελέσματα για όλες τις πειραματικές συνθήκες, συμπεριλαμβανομένων συνοπτικών στατιστικών και συγκρίσεων μεταξύ συνθηκών.

\section{Συνθήκη 1: Baseline Features (47) + Unweighted}
% Τα αποτελέσματα παράχθηκαν από το experimental pipeline (βλ. Chapter~\ref{ch:experimental-design})

\subsection{Dataset A: MDVR-KCL Σύνοψη}

\begin{table}[H]
\centering
\small
\begin{tabular}{@{}llccccc@{}}
\toprule
\textbf{Task} & \textbf{Model} & \textbf{Accuracy} & \textbf{Precision} & \textbf{Recall} & \textbf{F1} & \textbf{ROC-AUC} \\
\midrule
\multirow{3}{*}{ReadText} 
& LR & $0.621 \pm 0.058$ & $0.620 \pm 0.217$ & $0.550 \pm 0.201$ & $0.542 \pm 0.099$ & $0.717 \pm 0.139$ \\
& SVM & $0.621 \pm 0.106$ & $0.367 \pm 0.342$ & $0.317 \pm 0.335$ & $0.333 \pm 0.333$ & $0.614 \pm 0.312$ \\
& RF & $0.629 \pm 0.178$ & $0.450 \pm 0.447$ & $0.333 \pm 0.408$ & $0.351 \pm 0.363$ & $0.590 \pm 0.302$ \\
\midrule
\multirow{3}{*}{Spontaneous}
& LR & $0.639 \pm 0.160$ & $0.469 \pm 0.292$ & $0.667 \pm 0.408$ & $0.539 \pm 0.321$ & $0.760 \pm 0.214$ \\
& SVM & $0.636 \pm 0.135$ & $0.600 \pm 0.435$ & $0.333 \pm 0.236$ & $0.400 \pm 0.253$ & $0.407 \pm 0.309$ \\
& RF & $0.721 \pm 0.176$ & $0.633 \pm 0.415$ & $0.600 \pm 0.435$ & $0.567 \pm 0.365$ & $\mathbf{0.828 \pm 0.148}$ \\
\bottomrule
\end{tabular}
\caption{Αποτελέσματα Condition 1 ανά task (Dataset A, baseline features, χωρίς weighting)}
\label{tab:c1-summary}
\end{table}

\subsection{Dataset B: PD Speech Features Σύνοψη}

\begin{table}[H]
\centering
\small
\begin{tabular}{@{}lccccc@{}}
\toprule
\textbf{Model} & \textbf{Accuracy} & \textbf{Precision} & \textbf{Recall} & \textbf{F1} & \textbf{ROC-AUC} \\
\midrule
LogisticRegression & $0.828 \pm 0.008$ & $0.881 \pm 0.010$ & $0.890 \pm 0.016$ & $0.885 \pm 0.006$ & $0.867 \pm 0.029$ \\
SVM\_RBF & $0.851 \pm 0.024$ & $0.841 \pm 0.015$ & $0.986 \pm 0.019$ & $0.908 \pm 0.015$ & $0.885 \pm 0.025$ \\
RandomForest & $0.882 \pm 0.019$ & $0.876 \pm 0.012$ & $0.980 \pm 0.015$ & $0.925 \pm 0.012$ & $\mathbf{0.940 \pm 0.013}$ \\
\bottomrule
\end{tabular}
\caption{Αποτελέσματα Condition 1 (Dataset B, baseline features, χωρίς weighting)}
\label{tab:c1-datasetb}
\end{table}

\textbf{Σημείωση:} Το Dataset B εμφανίζει ουσιαστικά υψηλότερη απόδοση και χαμηλότερη διακύμανση από το Dataset A, εύρημα συμβατό με πιθανό optimistic bias λόγω άγνωστης επικάλυψης subjects.

\section{Συνθήκη 2: Extended Features (78) + Unweighted}
% Τα αποτελέσματα παράχθηκαν από το experimental pipeline (βλ. Chapter~\ref{ch:experimental-design})

\subsection{Dataset A: MDVR-KCL Σύνοψη}

\begin{table}[H]
\centering
\small
\begin{tabular}{@{}llccccc@{}}
\toprule
\textbf{Task} & \textbf{Model} & \textbf{Accuracy} & \textbf{Precision} & \textbf{Recall} & \textbf{F1} & \textbf{ROC-AUC} \\
\midrule
\multirow{3}{*}{ReadText} 
& LR & $0.596 \pm 0.079$ & $0.600 \pm 0.253$ & $0.433 \pm 0.149$ & $0.475 \pm 0.106$ & $0.698 \pm 0.132$ \\
& SVM & $0.786 \pm 0.181$ & $0.733 \pm 0.435$ & $0.567 \pm 0.365$ & $0.634 \pm 0.386$ & $0.834 \pm 0.153$ \\
& RF & $0.818 \pm 0.140$ & $0.883 \pm 0.162$ & $0.700 \pm 0.298$ & $0.746 \pm 0.207$ & $\mathbf{0.822 \pm 0.166}$ \\
\midrule
\multirow{3}{*}{Spontaneous}
& LR & $0.671 \pm 0.199$ & $0.500 \pm 0.361$ & $0.600 \pm 0.435$ & $0.530 \pm 0.377$ & $0.783 \pm 0.139$ \\
& SVM & $0.636 \pm 0.089$ & $0.533 \pm 0.361$ & $0.400 \pm 0.279$ & $0.428 \pm 0.258$ & $0.460 \pm 0.294$ \\
& RF & $0.779 \pm 0.161$ & $0.683 \pm 0.410$ & $0.600 \pm 0.435$ & $0.605 \pm 0.387$ & $\mathbf{0.857 \pm 0.171}$ \\
\bottomrule
\end{tabular}
\caption{Αποτελέσματα Condition 2 ανά task (Dataset A, extended features, χωρίς weighting)}
\label{tab:c2-summary}
\end{table}

\subsection{Βελτίωση έναντι Baseline (Random Forest, Dataset A)}

\begin{table}[H]
\centering
\begin{tabular}{@{}lccc@{}}
\toprule
\textbf{Task} & \textbf{$\Delta$ Accuracy} & \textbf{$\Delta$ ROC-AUC} & \textbf{$\Delta$ Variance} \\
\midrule
ReadText & +18.9pp & \textbf{+23.2pp} & Reduced (0.302 $\rightarrow$ 0.166) \\
SpontaneousDialogue & +5.8pp & +2.9pp & Increased (0.148 $\rightarrow$ 0.171) \\
\bottomrule
\end{tabular}
\caption{Βελτίωση Condition 2 έναντι Condition 1 (Random Forest)}
\label{tab:c2-improvement}
\end{table}

\textbf{Παρατήρηση:} Το feature extension βελτίωσε αισθητά την απόδοση στο ReadText (+23.2pp ROC-AUC), ενώ στο SpontaneousDialogue παρατηρήθηκαν πιο περιορισμένα gains (+2.9pp). Αυτό υποδηλώνει ότι τα extended features αποτυπώνουν αποτελεσματικότερα τις δυναμικές της δομημένης ομιλίας σε σύγκριση με την αυθόρμητη ομιλία.

\section{Συνθήκη 3: Baseline Features (47) + Weighted}
% Τα αποτελέσματα παράχθηκαν από το experimental pipeline (βλ. Chapter~\ref{ch:experimental-design})

\subsection{Dataset A: MDVR-KCL Σύνοψη}

\begin{table}[H]
\centering
\small
\begin{tabular}{@{}llccccc@{}}
\toprule
\textbf{Task} & \textbf{Model} & \textbf{Accuracy} & \textbf{Precision} & \textbf{Recall} & \textbf{F1} & \textbf{ROC-AUC} \\
\midrule
\multirow{3}{*}{ReadText} 
& LR & $0.596 \pm 0.079$ & $0.600 \pm 0.235$ & $0.550 \pm 0.201$ & $0.528 \pm 0.099$ & $0.717 \pm 0.139$ \\
& SVM & $0.704 \pm 0.111$ & $0.617 \pm 0.371$ & $0.500 \pm 0.373$ & $0.519 \pm 0.320$ & $0.542 \pm 0.312$ \\
& RF & $0.650 \pm 0.148$ & $0.550 \pm 0.371$ & $0.400 \pm 0.365$ & $0.431 \pm 0.306$ & $0.687 \pm 0.258$ \\
\midrule
\multirow{3}{*}{Spontaneous}
& LR & $0.639 \pm 0.160$ & $0.469 \pm 0.292$ & $0.667 \pm 0.408$ & $0.539 \pm 0.321$ & $0.760 \pm 0.214$ \\
& SVM & $0.693 \pm 0.123$ & $0.567 \pm 0.365$ & $0.600 \pm 0.365$ & $0.560 \pm 0.318$ & $0.423 \pm 0.312$ \\
& RF & $0.693 \pm 0.123$ & $0.583 \pm 0.373$ & $0.600 \pm 0.435$ & $0.538 \pm 0.326$ & $\mathbf{0.827 \pm 0.133}$ \\
\bottomrule
\end{tabular}
\caption{Αποτελέσματα Condition 3 ανά task (Dataset A, baseline features, με weighting)}
\label{tab:c3-summary}
\end{table}

\subsection{Επίδραση του Weighting (vs Condition 1, Random Forest)}

\begin{table}[H]
\centering
\begin{tabular}{@{}lccc@{}}
\toprule
\textbf{Task} & \textbf{$\Delta$ Accuracy} & \textbf{$\Delta$ ROC-AUC} & \textbf{Effect} \\
\midrule
ReadText & +2.1pp & +9.7pp & Improved \\
SpontaneousDialogue & $-2.8$pp & $-0.1$pp & Negligible \\
\bottomrule
\end{tabular}
\caption{Επίδραση του weighting στο Condition 3 έναντι Condition 1 (RF)}
\label{tab:c3-weighting}
\end{table}

\textbf{Παρατήρηση:} Το class weighting βελτίωσε την απόδοση στο ReadText, αλλά είχε αμελητέα επίδραση στο SpontaneousDialogue, γεγονός που υποδηλώνει task-dependent ευαισθησία στον χειρισμό του class imbalance.

\section{Συνθήκη 4: Extended Features (78) + Weighted}
% Τα αποτελέσματα παράχθηκαν από το experimental pipeline (βλ. Chapter~\ref{ch:experimental-design})

\subsection{Dataset A: MDVR-KCL Σύνοψη}

\begin{table}[H]
\centering
\small
\begin{tabular}{@{}llccccc@{}}
\toprule
\textbf{Task} & \textbf{Model} & \textbf{Accuracy} & \textbf{Precision} & \textbf{Recall} & \textbf{F1} & \textbf{ROC-AUC} \\
\midrule
\multirow{3}{*}{ReadText} 
& LR & $0.650 \pm 0.108$ & $0.650 \pm 0.253$ & $0.550 \pm 0.201$ & $0.564 \pm 0.160$ & $0.698 \pm 0.132$ \\
& SVM & $0.761 \pm 0.214$ & $0.700 \pm 0.447$ & $0.567 \pm 0.365$ & $0.620 \pm 0.390$ & $0.834 \pm 0.153$ \\
& RF & $0.818 \pm 0.140$ & $0.883 \pm 0.162$ & $0.700 \pm 0.298$ & $0.746 \pm 0.207$ & $0.805 \pm 0.182$ \\
\midrule
\multirow{3}{*}{Spontaneous}
& LR & $0.696 \pm 0.179$ & $0.520 \pm 0.356$ & $0.667 \pm 0.471$ & $0.563 \pm 0.381$ & $0.783 \pm 0.139$ \\
& SVM & $0.664 \pm 0.167$ & $0.567 \pm 0.365$ & $0.600 \pm 0.365$ & $0.560 \pm 0.318$ & $0.403 \pm 0.347$ \\
& RF & $0.721 \pm 0.203$ & $0.653 \pm 0.409$ & $0.600 \pm 0.435$ & $0.583 \pm 0.373$ & $\mathbf{0.823 \pm 0.209}$ \\
\bottomrule
\end{tabular}
\caption{Αποτελέσματα Condition 4 ανά task (Dataset A, extended features, με weighting)}
\label{tab:c4-summary}
\end{table}

\subsection{Επίδραση του Weighting (vs Condition 2, Random Forest)}

\begin{table}[H]
\centering
\begin{tabular}{@{}lccc@{}}
\toprule
\textbf{Task} & \textbf{$\Delta$ Accuracy} & \textbf{$\Delta$ ROC-AUC} & \textbf{Effect} \\
\midrule
ReadText & 0.0pp & $-1.7$pp & Slight decrease \\
SpontaneousDialogue & $-5.8$pp & $-3.4$pp & Decreased \\
\bottomrule
\end{tabular}
\caption{Επίδραση του weighting στο Condition 4 έναντι Condition 2 (RF)}
\label{tab:c4-weighting}
\end{table}

\textbf{Παρατήρηση:} Ο συνδυασμός extended features με class weighting δεν βελτίωσε την απόδοση σε σχέση με τη χρήση extended features μόνο. Αυτό υποδηλώνει ότι feature extension και class weighting αντιμετωπίζουν παρόμοιες πτυχές του classification problem, με diminishing returns όταν εφαρμόζονται από κοινού.

\section{Matrix Σύγκρισης μεταξύ Conditions}

\subsection{Random Forest ROC-AUC ανά Task}

\begin{table}[H]
\centering
\begin{tabular}{@{}llcc@{}}
\toprule
\textbf{Task} & & \textbf{Baseline (47)} & \textbf{Extended (78)} \\
\midrule
\multirow{2}{*}{ReadText}
& Unweighted & $0.590 \pm 0.302$ & $\mathbf{0.822 \pm 0.166}$ \\
& Weighted & $0.687 \pm 0.258$ & $0.805 \pm 0.182$ \\
\midrule
\multirow{2}{*}{Spontaneous}
& Unweighted & $0.828 \pm 0.148$ & $\mathbf{0.857 \pm 0.171}$ \\
& Weighted & $0.827 \pm 0.133$ & $0.823 \pm 0.209$ \\
\bottomrule
\end{tabular}
\caption{Πίνακας σύγκρισης Random Forest ROC-AUC (Dataset A)}
\label{tab:rf-rocauc-matrix}
\end{table}

\subsection{Random Forest Accuracy ανά Task}

\begin{table}[H]
\centering
\begin{tabular}{@{}llcc@{}}
\toprule
\textbf{Task} & & \textbf{Baseline (47)} & \textbf{Extended (78)} \\
\midrule
\multirow{2}{*}{ReadText}
& Unweighted & $62.9\% \pm 17.8\%$ & $\mathbf{81.8\% \pm 14.0\%}$ \\
& Weighted & $65.0\% \pm 14.8\%$ & $81.8\% \pm 14.0\%$ \\
\midrule
\multirow{2}{*}{Spontaneous}
& Unweighted & $72.1\% \pm 17.6\%$ & $\mathbf{77.9\% \pm 16.1\%}$ \\
& Weighted & $69.3\% \pm 12.3\%$ & $72.1\% \pm 20.3\%$ \\
\bottomrule
\end{tabular}
\caption{Πίνακας σύγκρισης Random Forest Accuracy (Dataset A)}
\label{tab:rf-accuracy-matrix}
\end{table}

\subsection{Κύριες Παρατηρήσεις}

\begin{enumerate}
    \item \textbf{Το feature extension βελτιώνει σταθερά το ReadText:} +23.2pp ROC-AUC (χωρίς weighting)
    \item \textbf{Το SpontaneousDialogue έχει υψηλότερο baseline:} 0.828 vs 0.590 ROC-AUC με baseline features
    \item \textbf{Το class weighting εμφανίζει μικτά αποτελέσματα:} βελτιώνει τη baseline επίδοση στο ReadText αλλά μειώνει την επίδοση με extended features
    \item \textbf{Υψηλή διακύμανση σε όλα τα settings:} τα standard deviations συχνά υπερβαίνουν το 0.15, αντανακλώντας το μικρό sample size ($n=37$)
\end{enumerate}

\section{Σημειώσεις για τη Statistical Significance}

Λόγω υψηλών standard deviations (συχνά $> 0.15$), τα confidence intervals επικαλύπτονται σε πολλές συγκρίσεις. Αυτό περιορίζει τη δυνατότητα για ισχυρούς στατιστικούς ισχυρισμούς σχετικά με τις διαφορές μεταξύ conditions. Τα αποτελέσματα θα πρέπει να ερμηνεύονται ως \textbf{trends} και όχι ως οριστικές κατατάξεις.

\section{Καμπύλες ROC}

Παρουσιάζονται καμπύλες Receiver Operating Characteristic (ROC) για όλους τους πέντε classifiers σε κάθε experimental condition. Οι αχνές καμπύλες αντιστοιχούν σε μεμονωμένα cross-validation folds, ενώ η έντονη καμπύλη αντιστοιχεί στον μέσο όρο των folds με shaded band $\pm 1$ standard deviation. Όλες οι καμπύλες δημιουργήθηκαν με χρήση του extended feature set.

\begin{figure}[H]
    \centering
    \includegraphics[width=\textwidth]{fig_roc_readtext.pdf}
    \caption{ROC curves για όλους τους πέντε classifiers στο ReadText task (Dataset A, extended features). Τα shaded bands υποδεικνύουν $\pm 1$ standard deviation στα πέντε cross-validation folds.}
    \label{fig:roc-readtext}
\end{figure}

\begin{figure}[H]
    \centering
    \includegraphics[width=\textwidth]{fig_roc_spontaneous.pdf}
    \caption{ROC curves για όλους τους πέντε classifiers στο SpontaneousDialogue task (Dataset A, extended features).}
    \label{fig:roc-spontaneous}
\end{figure}

\begin{figure}[H]
    \centering
    \includegraphics[width=\textwidth]{fig_roc_dataset_b.pdf}
    \caption{ROC curves για όλους τους πέντε classifiers στο Dataset B (PD Speech Features). Σημειώνεται ότι subject-level deduplication δεν είναι εφικτό για αυτό το dataset· τα αποτελέσματα ενδέχεται να είναι optimistic λόγω πιθανής επικάλυψης subjects μεταξύ folds.}
    \label{fig:roc-datasetb}
\end{figure}
